\SourcePage{295}{298}
\section{Second quantization: Bose statistics}

In the theory of systems containing large numbers of identical particles, a special method known as \emph{second quantization} is widely used. It is particularly indispensable in relativistic theory, where

\SourcePage{296}{299}
one must deal with systems in which the number of particles itself is variable\footnote{The second-quantization method was developed by Dirac for photons in the theory of radiation (1927), and was subsequently extended to fermions by Wigner and Jordan (E.~Wigner, P.~Jordan, 1928).}.

Let $\psi_1(\xi)$, $\psi_2(\xi)$,~\dots be some complete orthonormal set of one-particle stationary-state wave functions\footnote{As in §~61, $\xi$ denotes the collection comprising the particle coordinates and its spin projection $\sigma$; integration over $d\xi$ is understood to mean integration over the coordinates together with summation over $\sigma$.}. These may be particle states in an arbitrarily chosen external field, but usually one simply takes plane waves---the wave functions of a free particle with specified values of momentum (and spin projection). To make the spectrum of states discrete, the motion of the particles is then considered in a large but bounded region of space. For motion in a bounded volume, the eigenvalues of the momentum components form discrete sequences; the intervals between adjacent values are inversely proportional to the linear dimensions of the region and tend to zero as those dimensions increase.

In a system of free particles, the momentum of each particle is separately conserved. Consequently the \emph{occupation numbers} are conserved as well: the numbers $N_1,N_2,\ldots$ specifying how many particles occupy each of the states $\psi_1,\psi_2,\ldots$. In a system of interacting particles, the individual particle momenta are no longer conserved, and neither are the occupation numbers. For such a system one can speak only of the probability distribution of the various occupation-number values. Our aim is to construct a mathematical formalism in which the occupation numbers themselves, rather than particle coordinates and spin projections, serve as the independent variables.

Dirac notation (see the end of §~11) is convenient in this formalism, with $N_1,N_2,\ldots$ chosen as the quantum numbers specifying a state. The states corresponding to the wave functions (61.3) and (61.5) will be denoted by $|N_1,N_2,\ldots\rangle$. Coordinate and spin variables no longer appear explicitly.

With this choice of independent variables, the operators of the various physical quantities, including the Hamiltonian of the system, must likewise be formulated through their action on functions of the occupation numbers. Such a formulation can be reached by starting from the ordinary matrix representation of operators. For this purpose we must consider

\SourcePage{297}{300}
operator matrix elements between stationary-state wave functions of a system of noninteracting particles. Since these states can be described by prescribed occupation-number values, the way in which the operators act on these variables will thereby become clear.

We first consider systems of particles obeying Bose statistics.

Let $\hat f_a^{(1)}$ be the operator of some quantity associated with one particle (particle $a$), so that it acts only on functions of the variables $\xi_a$. Introduce the operator, symmetric in all the particles,
\begin{equation}
 \hat F^{(1)}=\sum_a\hat f_a^{(1)}
 \tag{64.1}
\end{equation}
(the sum extends over all particles), and determine its matrix elements between the wave functions (61.3). It is readily seen first of all that the matrix elements can be nonzero only for transitions in which the numbers $N_1,N_2,\ldots$ remain unchanged (the diagonal elements), and for transitions in which one of these numbers increases by unity while another decreases by unity. Indeed, every operator $\hat f_a^{(1)}$ acts on only one function in the product $\psi_{p_1}(\xi_1)\psi_{p_2}(\xi_2)\cdots\psi_{p_N}(\xi_N)$, so its matrix elements can be nonzero only for transitions in which the state of one particle changes. This means that the number of particles in one state decreases by unity and the number in another increases by unity. The actual calculation of these matrix elements is quite elementary; it is easier to carry it out than to follow a written account of it. We therefore give only the result. The off-diagonal elements are
\begin{equation}
 \langle N_i,N_k-1|\hat F^{(1)}|N_i-1,N_k\rangle
 =f_{ik}^{(1)}\sqrt{N_iN_k}.
 \tag{64.2}
\end{equation}
Only the indices in which the matrix element is not diagonal have been displayed; the others are suppressed for brevity. Here $f_{ik}^{(1)}$ is the matrix element
\begin{equation}
 f_{ik}^{(1)}=\int\psi_i^*(\xi)\hat f^{(1)}\psi_k(\xi)\,d\xi;
 \tag{64.3}
\end{equation}
since the operators $\hat f_a^{(1)}$ differ only in the labels of the variables on which they act, the integrals (64.3) do not depend on the index $a$, and that index has been omitted. The diagonal matrix elements of $\hat F^{(1)}$ are the mean values

\SourcePage{298}{301}
of $F^{(1)}$ in the states $\Psi_{N_1N_2\ldots}$. Their calculation gives
\begin{equation}
 \overline{F^{(1)}}=\sum_i f_{ii}^{(1)}N_i.
 \tag{64.4}
\end{equation}

We now introduce the operators fundamental to second quantization, $\hat a_i$, which act not on functions of coordinates but on functions of the occupation numbers. By definition, when $\hat a_i$ acts on the state $|N_1,N_2,\ldots\rangle$, it lowers the variable $N_i$ by unity and at the same time multiplies the function by $\sqrt{N_i}$\footnote{The natural notation $\hat a|n\rangle$ has been used for the result of applying the operator $\hat a$ to the wave function of the state $|n\rangle$.}:
\begin{equation}
 \hat a_i|N_1,N_2,\ldots,N_i,\ldots\rangle
 =\sqrt{N_i}|N_1,N_2,\ldots,N_i-1,\ldots\rangle.
 \tag{64.5}
\end{equation}
Thus $\hat a_i$ may be said to reduce by unity the number of particles in state $i$; it is therefore called the particle \emph{annihilation operator}. It can be represented by a matrix whose sole nonzero element is
\begin{equation}
 \langle N_i-1|\hat a_i|N_i\rangle=\sqrt{N_i}.
 \tag{64.6}
\end{equation}

The operator $\hat a_i^+$ adjoint to $\hat a_i$ is represented, by definition (see (11.9)), by a matrix with the single element
\begin{equation}
 \langle N_i|\hat a_i^+|N_i-1\rangle
 =\langle N_i-1|\hat a_i|N_i\rangle^*=\sqrt{N_i}.
 \tag{64.7}
\end{equation}
It follows that, when this operator acts on $|N_1,N_2,\ldots\rangle$, it increases $N_i$ by 1:
\begin{equation}
 \hat a_i^+|N_1,N_2,\ldots,N_i,\ldots\rangle
 =\sqrt{N_i+1}|N_1,N_2,\ldots,N_i+1,\ldots\rangle.
 \tag{64.8}
\end{equation}
In other words, $\hat a_i^+$ increases by 1 the number of particles in state $i$; it is called the particle \emph{creation operator}.

When the product $\hat a_i^+\hat a_i$ acts on a wave function, it can only multiply it by a constant, leaving all the variables $N_1,N_2,\ldots$ unchanged: $\hat a_i$ first lowers $N_i$ by 1, after which $\hat a_i^+$ restores its original value. Direct multiplication of the matrices (64.6) and (64.7) indeed shows that $\hat a_i^+\hat a_i$ is represented by a diagonal matrix whose diagonal elements are $N_i$:
\begin{equation}
 \hat a_i^+\hat a_i=N_i.
 \tag{64.9}
\end{equation}
Similarly,
\begin{equation}
 \hat a_i\hat a_i^+=N_i+1.
 \tag{64.10}
\end{equation}

\SourcePage{299}{302}
Subtracting these expressions gives the commutation rule for $\hat a_i$ and $\hat a_i^+$:
\begin{equation}
 \hat a_i\hat a_i^+-\hat a_i^+\hat a_i=1.
 \tag{64.11}
\end{equation}
Operators with different indices $i$ and $k$, which act on different variables ($N_i$ and $N_k$), commute:
\begin{equation}
 \hat a_i\hat a_k-\hat a_k\hat a_i=0,\qquad
 \hat a_i\hat a_k^+-\hat a_k^+\hat a_i=0,\qquad i\ne k.
 \tag{64.12}
\end{equation}

From the properties of $\hat a_i$, $\hat a_i^+$ just described, it is easy to see that the operator
\begin{equation}
 \hat F^{(1)}=\sum_{i,k}f_{ik}^{(1)}\hat a_i^+\hat a_k
 \tag{64.13}
\end{equation}
is identical to the operator (64.1). Indeed, all the matrix elements calculated using (64.6) and (64.7) coincide with (64.2) and (64.4). This is a very important result. In (64.13) the quantities $f_{ik}^{(1)}$ are simply numbers. We have therefore succeeded in expressing an ordinary operator acting on functions of coordinates as an operator acting on functions of the new variables, the occupation numbers $N_i$.

The result is easily generalized to operators of other forms. Let
\begin{equation}
 \hat F^{(2)}=\sum_{a>b}\hat f_{ab}^{(2)},
 \tag{64.14}
\end{equation}
where $\hat f_{ab}^{(2)}$ is the operator of a physical quantity associated with a pair of particles at once, and therefore acts on functions of $\xi_a$ and $\xi_b$. An analogous calculation shows that such an operator can be expressed in terms of $\hat a_i$, $\hat a_i^+$ as
\begin{equation}
 \hat F^{(2)}=\frac12\sum_{i,k,l,m}
 \langle ik|f^{(2)}|lm\rangle\hat a_i^+\hat a_k^+\hat a_m\hat a_l,
 \tag{64.15}
\end{equation}
\[
 \langle ik|f^{(2)}|lm\rangle=\iint
 \psi_i^*(\xi_1)\psi_k^*(\xi_2)\hat f^{(2)}
 \psi_l(\xi_1)\psi_m(\xi_2)\,d\xi_1d\xi_2.
\]
The extension of these formulas to operators of any other form symmetric in all particles ($\hat F^{(3)}=\sum\hat f_{abc}^{(3)}$, and so on) is evident.

The same formulas can be used to express in terms of $\hat a_i$, $\hat a_i^+$ the Hamiltonian of the physical system under consideration, consisting of $N$ interacting identical particles. The Hamiltonian of such a system is, of course, symmetric in all the particles. In the nonrelativistic approximation\footnote{In the absence of a magnetic field.} it is independent of the particle spins

\SourcePage{300}{303}
and can be written in the general form
\begin{equation}
 \hat H=\sum_a\hat H_a^{(1)}+\sum_{a>b}U^{(2)}(\mathbf r_a,\mathbf r_b)
 +\sum_{a>b>c}U^{(3)}(\mathbf r_a,\mathbf r_b,\mathbf r_c)+\cdots
 \tag{64.16}
\end{equation}
Here $\hat H_a^{(1)}$ is the part of the Hamiltonian that depends on the coordinates of only one particle, particle $a$:
\begin{equation}
 \hat H_a^{(1)}=-\frac{\hbar^2}{2m}\Delta_a+U^{(1)}(\mathbf r_a),
 \tag{64.17}
\end{equation}
where $U^{(1)}(\mathbf r_a)$ is the potential energy of one particle in the external field. The remaining terms in (64.16) represent the interaction energy of the particles with one another, with terms depending respectively on the coordinates of two, three, and so forth particles separated out.

Writing the Hamiltonian in this form permits direct application of (64.13), (64.15), and their analogues. Thus,
\begin{equation}
 \hat H=\sum_{i,k}H_{ik}^{(1)}\hat a_i^+\hat a_k
 +\frac12\sum_{i,k,l,m}\langle ik|U^{(2)}|lm\rangle
 \hat a_i^+\hat a_k^+\hat a_m\hat a_l+\cdots
 \tag{64.18}
\end{equation}
This is the required expression of the Hamiltonian as an operator acting on functions of the occupation numbers.

For a system of noninteracting particles, only the first term remains in (64.18):
\begin{equation}
 \hat H=\sum_{i,k}H_{ik}^{(1)}\hat a_i^+\hat a_k.
 \tag{64.19}
\end{equation}
If the functions $\psi_i$ are chosen as eigenfunctions of the one-particle Hamiltonian $\hat H^{(1)}$, the matrix $H_{ik}^{(1)}$ is diagonal and its diagonal elements are the particle energy eigenvalues $\varepsilon_i$. Hence
\[
 \hat H=\sum_i\varepsilon_i\hat a_i^+\hat a_i;
\]
replacing the operator $\hat a_i^+\hat a_i$ by its eigenvalues (64.9), we obtain the energy levels of the system in the form
\[
 E=\sum_i\varepsilon_iN_i,
\]
the trivial result that was to be expected.

\SourcePage{301}{304}
The formalism developed here may be put in a more compact form by introducing the so-called $\psi$-operators\footnote{Notice the analogy between (64.20) and the expansion $\psi=\sum a_i\psi_i$ of an arbitrary wave function in a complete set of functions. Here the expansion is, as it were, quantized once again, which is the origin of the name of the method: second quantization.}
\begin{equation}
 \hat\psi(\xi)=\sum_i\psi_i(\xi)\hat a_i,\qquad
 \hat\psi^+(\xi)=\sum_i\psi_i^*(\xi)\hat a_i^+,
 \tag{64.20}
\end{equation}
where the variables $\xi$ are treated as parameters. From the properties of $\hat a_i$, $\hat a_i^+$ stated above, it is clear that $\hat\psi$ lowers the total number of particles in the system by unity, while $\hat\psi^+$ raises it by unity.

It is easy to see that $\hat\psi^+(\xi_0)$ creates a particle located at the point $\xi_0$. Indeed, the action of $\hat a_i^+$ creates a particle in the state with wave function $\psi_i(\xi)$. Consequently, the action of $\psi_i^*(\xi_0)\hat a_i^+$ creates a particle in a state whose wave function is
\[
 \sum_i\psi_i^*(\xi)\psi_i(\xi_0)=\delta(\xi-\xi_0)
\]
(formula (5.12) has been used), which corresponds to a particle with definite coordinate and spin values\footnote{$\delta(\xi-\xi_0)$ is a shorthand for the product
\[
 \delta(x-x_0)\delta(y-y_0)\delta(z-z_0)\delta_{\sigma\sigma_0}.
\]}).

The commutation rules for the $\psi$-operators follow directly from those for $\hat a_i$, $\hat a_i^+$:
\begin{equation}
 \hat\psi(\xi)\hat\psi(\xi')-\hat\psi(\xi')\hat\psi(\xi)=0,
 \tag{64.21}
\end{equation}
\begin{equation}
 \hat\psi(\xi)\hat\psi^+(\xi')-\hat\psi^+(\xi')\hat\psi(\xi)
 =\sum_i\psi_i(\xi)\psi_i^*(\xi')=\delta(\xi-\xi').
 \tag{64.22}
\end{equation}

In terms of the $\psi$-operators, the second-quantized operator $\hat F^{(1)}$ is
\begin{equation}
 \hat F^{(1)}=\int\hat\psi^+(\xi)\hat f^{(1)}\hat\psi(\xi)\,d\xi
 \tag{64.23}
\end{equation}
(the operator $\hat f^{(1)}$ is understood here to act in $\hat\psi(\xi)$ on functions of the parameters $\xi$). Indeed, inserting $\hat\psi$ and $\hat\psi^+$ from (64.20) and using definition (64.3), we recover

\SourcePage{302}{305}
formula (64.13). Similarly, in place of (64.15) we have
\begin{equation}
 \hat F^{(2)}=\frac12\iint\hat\psi^+(\xi)\hat\psi^+(\xi')
 \hat f^{(2)}\hat\psi(\xi')\hat\psi(\xi)\,d\xi\,d\xi'.
 \tag{64.24}
\end{equation}

In particular, the Hamiltonian of the system expressed through the $\psi$-operators takes the form
\begin{equation}
 \begin{aligned}
 \hat H={}&\int\left\{-\frac{\hbar^2}{2m}\hat\psi^+(\xi)
 \Delta\hat\psi(\xi)+\hat\psi^+(\xi)U^{(1)}(\xi)\hat\psi(\xi)\right\}d\xi\\
 &+\frac12\iint\hat\psi^+(\xi)\hat\psi^+(\xi')U^{(2)}(\xi,\xi')
 \hat\psi(\xi')\hat\psi(\xi)\,d\xi\,d\xi'+\cdots
 \end{aligned}
 \tag{64.25}
\end{equation}

The operator $\hat\psi^+(\xi)\hat\psi(\xi)$, constructed from the $\psi$-operators in analogy with the product $\psi^*\psi$ that gives the probability density of a particle in a state with wave function $\psi$, is called the particle \emph{density operator}. The integral
\begin{equation}
 \hat N=\int\hat\psi^+\hat\psi\,d\xi
 \tag{64.26}
\end{equation}
plays the role in second quantization of the operator for the total number of particles in the system. Indeed, substituting the $\psi$-operators from (64.20) and using the normalization and mutual orthogonality of the wave functions gives $\hat N=\sum_i\hat a_i^+\hat a_i$. Every term in this sum is the particle-number operator for state $i$: according to (64.9), its eigenvalues are the occupation numbers $N_i$, and their sum is the total number of particles in the system\footnote{For systems with a prescribed number of particles, these statements (like the properties of the free-particle Hamiltonian (64.19)) appear trivial. Their extension in relativistic theory, however, leads to new and far from trivial results (cf. Vol.~IV, §~11).}.

Finally, if the system contains bosons of different kinds, second quantization requires a separate set of operators $\hat a$, $\hat a^+$ for every kind of particle. Operators belonging to different kinds of particles plainly commute with one another.
